% Einheit 1 von 4 – Erwartungswert berechnen und deuten: die Verteilung beschaffen (Tabelle, Sachtext, Baumpfade, Kosten je Ausgang; fehlende Wahrscheinlichkeit über die Summe 1), die gewichtete Summe bilden (bei Anzahlen n · p), das Ergebnis deuten (bank/kenngroessen-von-verteilungen/e1.jsonl, zusammenbau v0.1)
%% TODO Zeitmarke Einheit 1: Marken-Zeile nicht lesbar
%% TODO Zweigzeile Teil 1: was hier gelernt wird (Satz in Schülersprache aus dem Einheitstitel) – Einheit 1
\einheitenkopf[e1]{Einheit 1 von 4 · Erwartungswert berechnen und deuten: die Verteilung beschaffen (Tabelle, Sachtext, Baumpfade, Kosten je Ausgang; fehlende Wahrscheinlichkeit über die Summe 1), die gewichtete Summe bilden (bei Anzahlen n · p), das Ergebnis deuten}
\zweigzeile{Abitur GK}
\setcounter{aufgabe}{8}

%% TODO Titel: Ich-kann-Satz fehlt in der Bank (Platzhalter: Kettenname) – Nr. 9
%% TODO Anweisung: ein Satz über der ersten Teilaufgabe fehlt in der Bank – Nr. 9
\begin{aufgabe}{Erwartungswert berechnen und deuten}
%% TODO \rechenplatz{n}: Schrittzahl des Grundfalls fehlt in der Bank (nur bei fester Schreibform, 2.3 b)
\begin{teile}
\teil Ein Würfelspiel kostet 1\,€ Einsatz. $X$ ist der Betrag in Euro, den der Spieler nach dem Wurf erhält. Was beschreibt $X$? \\ \kreuz{Auszahlung} \\ \kreuz{Gewinn}
\teil Die Tabelle zeigt die Verteilung der Zufallsgröße $X$; eine Wahrscheinlichkeit fehlt. Ergänze sie und berechne den Erwartungswert $E(X)$.

\sachtabelle{lccc}{$x$ & 0 & 1 & 2}{$P(X = x)$ & 0{,}5 & 0{,}3 & \leerzelle}
fehlend: \leerfeld , $E(X) = $ \leerfeld
\teil Bei einem Losverkauf bringen 2\,\% der Lose 50\,€, 10\,\% der Lose 5\,€, die übrigen nichts. $X$ ist die Auszahlung für ein Los in Euro. Wie groß ist $E(X)$? $E(X) = $ \leerfeld[€]
\teil Ein Spiel kostet 3\,€ Einsatz. Mit der Wahrscheinlichkeit 0{,}25 werden 10\,€ ausgezahlt, sonst nichts. Zeige, dass das Spiel für den Anbieter auf lange Sicht gewinnbringend ist, und gib seinen mittleren Gewinn je Spiel an \hfill (Abitur 2018 GK). \leerfeld[€] je Spiel
\teil Eine Bäckerei verkauft Frühstückstüten: ein Brötchen für 1\,€, dazu Käse für 2\,€ oder Lachs für 4\,€ (beide gleich oft gewählt) und ein Kaffee für 2\,€. Wie hoch sind die Einnahmen bei 120 Tüten, und wie viele Tüten müssen mindestens verkauft werden, damit die Einnahmen mindestens 1\,000\,€ betragen \hfill (FHR 2021)? Einnahmen: \leerfeld[€] , mindestens \leerfeld Tüten
\teil In einer Urne liegen drei rote und sieben blaue Kugeln. Ein Spieler zieht zweimal mit Zurücklegen. Bei zwei roten Kugeln erhält er 9\,€, bei genau einer roten Kugel 4\,€, sonst nichts. Der Term $9 \cdot 0{,}3^2 + 4 \cdot 2 \cdot 0{,}3 \cdot 0{,}7 + 0 \cdot 0{,}7^2$ gibt den Erwartungswert der Auszahlung in Euro an. Erläutere die Bedeutung des zweiten Summanden im Sachzusammenhang \hfill (Abitur 2023 GK).
\end{teile}
\end{aufgabe}

%% TODO Titel: Ich-kann-Satz fehlt in der Bank (Platzhalter: Kettenname) – Nr. 10
%% TODO Anweisung: ein Satz über der ersten Teilaufgabe fehlt in der Bank – Nr. 10
\begin{aufgabe}{Prozentuale Abweichung einer Anzahl vom Erwartungswert berechnen}
\begin{teile}
\teil In einer Gärtnerei keimen erfahrungsgemäß 6\,\% der Samen nicht. Von 150 Samen keimen 12 nicht. Um wie viel Prozent weicht diese Anzahl vom Erwartungswert ab? Runde auf eine Stelle \hfill (Abitur 2022 GK). \leerfeld \%
\end{teile}
\end{aufgabe}

%% TODO Titel: Ich-kann-Satz fehlt in der Bank (Platzhalter: Kettenname) – Nr. 11
%% TODO Anweisung: ein Satz über der ersten Teilaufgabe fehlt in der Bank – Nr. 11
\begin{aufgabe}{Erwartungswert der Augensumme aus dem Erwartungswert der Trefferzahl berechnen}
\begin{teile}
\teil Ein Würfel trägt auf vier Seiten eine 1 und auf zwei Seiten eine 4. Er wird zwölfmal geworfen; $T$ ist die Anzahl der Vieren. Berechne $E(T)$ und daraus den Erwartungswert der Augensumme. $E(T) = $ \leerfeld , Augensumme: \leerfeld
\end{teile}
\end{aufgabe}

%% TODO Titel: Ich-kann-Satz fehlt in der Bank (Platzhalter: Kettenname) – Nr. 12
%% TODO Anweisung: ein Satz über der ersten Teilaufgabe fehlt in der Bank – Nr. 12
\begin{aufgabe}{Gesamtzahl der Gruppen aus der Personenzahl und der mittleren Gruppengröße unter einer vereinfachenden Annahme schätzen}
\begin{teile}
\teil Auf einem Campingplatz übernachten 2\,420 Personen in Zelten. Belegung der Zelte: 12\,\% mit einer Person, 46\,\% mit zwei, 30\,\% mit drei und 12\,\% mit mindestens vier Personen. Schätze die Zahl der Zelte und nenne deine Annahme \hfill (Abitur 2017 LK). etwa \leerfeld Zelte
\end{teile}
\end{aufgabe}

%% TODO Titel: Ich-kann-Satz fehlt in der Bank (Platzhalter: Kettenname) – Nr. 13
%% TODO Anweisung: ein Satz über der ersten Teilaufgabe fehlt in der Bank – Nr. 13
\begin{aufgabe}{Erwartungswert berechnen und deuten – Fehler finden, Begründen, Darstellungswechsel, Anwendung}
\begin{teile}
\teil Ein Spiel kostet 2\,€ Einsatz; mit der Wahrscheinlichkeit 0{,}3 werden 6\,€ ausgezahlt, sonst nichts. Ben rechnet so. Finde den Fehler und rechne richtig. \rechnung{E &= 6 \cdot 0{,}3 = 1{,}80 && \text{also gewinne ich im Mittel } 1{,}80\text{ € je Spiel}}
\teil Der Erwartungswert der Augenzahl eines Würfels ist 3{,}5. Begründe, warum ein Erwartungswert kein möglicher Wert der Zufallsgröße sein muss.
\teil Die Zufallsgröße $X$ hat die Verteilung $P(X = 0) = 0{,}1$, $P(X = 1) = 0{,}4$, $P(X = 2) = 0{,}3$, $P(X = 3) = 0{,}2$. Zeichne das Säulendiagramm der Verteilung.

\saeulen{0/,1/,2/,3/}{0.5}{0.1}{$P(X = x)$}
\teil Eine Fahrradversicherung kostet 60\,€ im Jahr. Erfahrungsgemäß werden 4\,\% der versicherten Räder im Jahr gestohlen (Erstattung 800\,€), bei 10\,\% gibt es einen Schaden, für den die Versicherung 150\,€ zahlt. Mit welchem Überschuss je Vertrag kann die Versicherung auf lange Sicht rechnen? \leerfeld[€] je Vertrag
\end{teile}
\end{aufgabe}
